\chapter{Comprendre le domaine}\label{ch:domaine}

\begin{objectifs}[Objectifs --- étape 1, \texttt{forge-domaine}]
\begin{itemize}
  \item mener ou restituer un Event Storming et repérer les points chauds ;
  \item écrire un glossaire où chaque terme a une définition et un nom technique uniques ;
  \item classer les sous-domaines et dessiner une context map ;
  \item faire respecter le vocabulaire par un script, pas par la bonne volonté.
\end{itemize}
\end{objectifs}

\begin{situation}[Sur le terrain : \sika{}]
Atelier d'une demi-journée avec la trésorière générale et trois trésorières de groupe. Au bout d'une
heure, un désaccord surgit : pour l'une, \q{le pot} est la somme collectée ; pour l'autre, c'est la
remise elle-même ; une troisième dit \q{la cagnotte}. Si le code nomme \cmd{Pot}, \cmd{Jackpot} et
\cmd{Fund} trois choses identiques, le bogue est déjà écrit.
\end{situation}

\section{Event Storming : voir le métier se dérouler}
\begin{concept}[Event Storming]
Un atelier où le métier et la technique posent sur un mur, dans l'ordre chronologique, les
\emph{événements} qui se produisent (au participe passé : \q{Cotisation confirmée}), puis les
commandes qui les déclenchent, les acteurs, les politiques (\q{dès que\dots, alors\dots}), les systèmes
externes et, en rouge, les \emph{points chauds}, c'est-à-dire les questions sans réponse.
\end{concept}

\noindent\begin{tabularx}{\linewidth}{lXX}
\toprule
\textbf{Couleur} & \textbf{Élément} & \textbf{Exemple Sika} \\ \midrule
Orange & Événement métier & Cotisation confirmée \\
Bleu & Commande & Cotiser pour le tour \\
Jaune & Acteur & Trésorière de groupe \\
Lilas & Politique & Dès que tous ont cotisé, clôturer le tour \\
Rose & Système externe & Plateforme Mobile Money \\
Rouge & Point chaud & Que faire d'un membre qui ne cotise pas ? \\
\bottomrule
\end{tabularx}

\begin{figure}[H]
\centering
\resizebox{\linewidth}{!}{%
\begin{tikzpicture}[
  ev/.style={fill=ochre!35,draw=ochre,rounded corners=2pt,align=center,font=\sffamily\scriptsize,text width=2.4cm,minimum height=9mm},
  cde/.style={fill=accentl!35,draw=accentl,rounded corners=2pt,align=center,font=\sffamily\scriptsize,text width=2.4cm,minimum height=9mm},
  po/.style={fill=violet!15,draw=violet!60,rounded corners=2pt,align=center,font=\sffamily\scriptsize,text width=2.4cm,minimum height=9mm},
  hot/.style={fill=brick!20,draw=brick,rounded corners=2pt,align=center,font=\sffamily\scriptsize,text width=3.6cm},
  ar/.style={-{Stealth},grey,thick}]
  \node[ev] (e0) at (0,0) {Tour ouvert};
  \node[cde] (c1) at (3.2,0) {Cotiser};
  \node[ev] (e1) at (6.4,0) {Cotisation initiée};
  \node[ev] (e2) at (9.6,0.9) {Cotisation confirmée};
  \node[ev] (e3) at (9.6,-0.9) {Cotisation échouée};
  \node[po] (p1) at (12.8,0.9) {Dès que tous ont cotisé : clôturer};
  \node[ev] (e4) at (16,0.9) {Tour clôturé};
  \node[cde] (c2) at (16,-0.9) {Remettre le pot};
  \node[ev] (e5) at (12.8,-0.9) {Pot remis};
  \draw[ar] (e0)--(c1); \draw[ar] (c1)--(e1); \draw[ar] (e1)--(e2); \draw[ar] (e1)--(e3);
  \draw[ar] (e2)--(p1); \draw[ar] (p1)--(e4); \draw[ar] (e4)--(c2); \draw[ar] (c2)--(e5);
  \node[hot] at (9.6,-2.7) {\textbf{H2} Au bout de combien de temps une cotisation sans réponse est-elle \q{échouée} ?};
  \node[hot] at (14.4,-2.7) {\textbf{H1} Que fait-on d'un membre qui ne cotise pas avant la date limite ?};
\end{tikzpicture}}
\caption{Le flux principal de \sika{}, tel qu'il sort de l'atelier. Les deux points chauds sont en rouge.}
\end{figure}

\begin{piege}[Piège : trancher un point chaud sans son responsable]
H1 touche au règlement intérieur ; seule la présidente peut y répondre. Le skill ne tranche \emph{jamais}
un point chaud lui-même : il l'inscrit dans la table avec un responsable et une échéance. H2, lui,
conditionnera la règle R2 de la spécification (chapitre~\ref{ch:spec}).
\end{piege}

\section{Le glossaire : un seul mot par chose}
\begin{concept}[Langage omniprésent]
Un vocabulaire unique, partagé par le métier et le code. \emph{Un terme} = \emph{une définition} =
\emph{un nom technique anglais}. Les synonymes courants sont listés pour être \textbf{bannis}.
\end{concept}

Voici le glossaire de \sika{} (extrait). La dernière colonne alimente un script.
\noindent\begin{tabularx}{\linewidth}{llX>{\raggedright\arraybackslash}p{3.4cm}}
\toprule
\textbf{Terme} & \textbf{Nom technique} & \textbf{Définition} & \textbf{Synonymes bannis} \\ \midrule
Membre & \cmd{Member} & Personne qui cotise dans un groupe & adhérente, user, client \\
Groupe & \cmd{SavingsGroup} & Ensemble de membres qui cotisent ensemble & cercle \\
Tour & \cmd{Round} & Période au bout de laquelle un seul membre reçoit le pot & cycle, session \\
Cotisation & \cmd{Contribution} & Somme versée par un membre pour un tour & versement, payment \\
Pot & \cmd{Pot} & Somme des cotisations confirmées d'un tour & cagnotte, caisse \\
Bénéficiaire & \cmd{Beneficiary} & Membre qui reçoit le pot d'un tour & gagnant, winner \\
\bottomrule
\end{tabularx}

\subsection*{Le garde-fou FF-05 en action}
Un développeur écrit, par habitude, une classe \cmd{Payment}. Le script lit la colonne des synonymes et cherche
chaque mot dans le code (sortie réelle) :
\begin{code}
\begin{Verbatim}
$ scripts/check-glossary-terms.sh docs/01-domaine/glossaire.md src/
src/a.ts:1:class Payment { amountXof: number }
KO: terme banni « payment »
\end{Verbatim}
\end{code}
Après correction en \cmd{Contribution} :
\begin{code}
\begin{Verbatim}
OK: aucun terme banni
\end{Verbatim}
\end{code}

\begin{piege}[Piège : affaiblir le glossaire pour faire passer le contrôle]
Quand le script crie, la mauvaise réponse est de retirer le synonyme de la liste. La bonne :
renommer dans le code. Le glossaire est un contrat avec le métier.
\end{piege}

\section{Sous-domaines et context map}
\begin{concept}[Bounded context]
Une frontière à l'intérieur de laquelle un modèle et son vocabulaire sont cohérents. Hors de la
frontière, le même mot peut signifier autre chose.
\end{concept}
Les sous-domaines se classent en trois familles : le \textbf{cœur} (ce qui différencie, là où va l'effort de
conception), le \textbf{support} (nécessaire, pas différenciant) et le \textbf{générique} (commun à tous les systèmes, on
l'achète ou on le fait au plus simple).

\begin{figure}[H]
\centering
\resizebox{.95\linewidth}{!}{%
\begin{tikzpicture}[
  bc/.style={draw,rounded corners=4pt,align=center,font=\sffamily\small,minimum width=3.4cm,minimum height=12mm},
  core/.style={bc,fill=accent,draw=accent,text=white},
  sup/.style={bc,fill=ochre!15,draw=ochre},
  gen/.style={bc,fill=grey!15,draw=grey},
  rel/.style={-{Stealth},thick,accent},
  lab/.style={font=\sffamily\scriptsize,text=accent,fill=white,inner sep=1.5pt}]
  \node[core] (c) at (0,0) {\textbf{Cotisation}\\{\scriptsize cœur : tours, cotisations, pot}};
  \node[sup] (n) at (5.6,1.8) {\textbf{Notifications}\\{\scriptsize support : SMS}};
  \node[gen] (i) at (-5.6,1.8) {\textbf{Identité}\\{\scriptsize générique}};
  \node[gen] (p) at (0,-3) {\textbf{Paiement mobile}\\{\scriptsize générique : agrégateur Mobile Money}};
  \draw[rel] (c)--(n) node[midway,lab]{Customer / Supplier};
  \draw[rel] (i)--(c) node[midway,lab]{Conformist};
  \draw[rel] (p)--(c) node[midway,lab,xshift=0pt]{ACL};
\end{tikzpicture}}
\caption{La context map de \sika{} : un cœur, un support, deux contextes génériques.}
\end{figure}

\noindent\begin{tabularx}{\linewidth}{>{\raggedright\arraybackslash}p{3.4cm}>{\raggedright\arraybackslash}p{3.8cm}X}
\toprule
\textbf{Pattern} & \textbf{Entre} & \textbf{Pourquoi} \\ \midrule
ACL (couche anticorruption) & Paiement mobile $\to$ Cotisation & isoler les particularités de chaque opérateur ; le cœur ne parle jamais MTN ou Moov directement \\
Conformist & Identité $\to$ Cotisation & on adopte le modèle de la solution d'identité sans le déformer \\
Customer / Supplier & Cotisation $\to$ Notifications & le cœur exprime ses besoins ; le support s'y conforme \\
\bottomrule
\end{tabularx}

\begin{vousjouez}[À vous de jouer]
Listez dix événements de votre projet, en quelques minutes, au participe passé. Repérez un mot qu'un collègue
et vous employez pour deux choses différentes (ou deux mots pour une chose). Écrivez la ligne de glossaire correspondante, avec
son nom technique anglais et deux synonymes à bannir.
\end{vousjouez}

\begin{retenir}[À retenir]
\begin{itemize}
  \item Un point chaud n'est jamais tranché sans son responsable.
  \item Un terme = une définition = un nom technique ; les synonymes sont bannis et contrôlés (FF-05).
  \item L'effort de conception va au cœur ; le reste s'achète ou reste simple.
\end{itemize}
\end{retenir}
\source{skills/forge-domaine/SKILL.md, scripts/check-glossary-terms.sh}

\begin{acquis}
  \item Pourquoi un événement s'écrit-il au participe passé ?
  \item Que contient la colonne « Synonymes à éviter », et qui la lit ?
  \item Dans \sika{}, pourquoi le paiement mobile est-il isolé derrière une ACL ?
\end{acquis}
